\section{Cutting into pants and varying twists}
\label{sec:twists}

\subsection{What data a cut curve retains}

Fix an oriented hyperbolic surface of finite type, without cone points,
and a pants decomposition with internal cuff lengths $a_1,\ldots,a_N>0$,
where $N=3g-3+n$. Boundary lengths, including cusps, are fixed.
Apart from components parallel to a cutting cuff, cutting a simple
multicurve gives disjoint essential arcs in each pair of pants.
Cutting a general immersed curve, even a filling one,
need not give disjoint arcs: self-intersections may remain inside a pair
of pants. Nor are the actual segments generally orthogonal to the cuffs.

For example, twisting a curve once around a cuff leaves the number
of its intersections with that cuff unchanged but changes the way
its endpoints are joined across the cut. Dehn--Thurston coordinates
record intersection counts together with integral twisting data
(and the usual conventions for components parallel to cuffs).
An itinerary is the more detailed record of successive pants,
relative arc classes and crossings followed by an immersed curve.

Boundary intersection counts alone lose the relative homotopy classes of
the arcs, their cyclic joining pattern, and winding information. For simple
multicurves this information is organized by Dehn--Thurston coordinates.
For a general immersed curve one must retain a fuller itinerary. Once this
topological data is fixed, varying the metric twist is a well-defined and
tractable problem.

There are two complementary exact descriptions. One uses arc lengths as
generating functions in endpoint positions. The other uses products of
holonomy matrices. Both keep the information that intersection counts alone
discard.

\subsection{An exact local distance function}

A seam is a shortest perpendicular arc between two cuffs.
In the universal covering plane the relevant cuff lifts are
\emph{ultraparallel}: disjoint geodesics with distinct ideal
endpoints and a unique common perpendicular. The variables below
allow the arc endpoints to move away from the feet of that
perpendicular, so they describe ordinary crossing segments as
well as seams.

Take two distinct ultraparallel geodesic lifts of positive-length cuffs
in $\HH^2$, at common perpendicular distance $d>0$.
Let $u,v$ be signed distances on the two lines
from the perpendicular feet, with compatible orientations. Then
\begin{equation}
 \boxed{
 D_d(u,v)=\arcosh\bigl(\cosh d\cosh u\cosh v-\sinh u\sinh v\bigr).}
 \label{eq:localdistance}
\end{equation}
For example, in the hyperboloid model use
\[
 p(u)=(\cosh u,\sinh u,0),\qquad
 q(v)=(\cosh d\cosh v,\sinh v,\sinh d\cosh v).
\]
Then $\cosh D=-\langle p,q\rangle$ for the Lorentz form of signature
$(-,+,+)$, and $\arcosh$ is its nonnegative real branch.
Reversing one endpoint coordinate changes the sign of the mixed term.

Setting $u=v=0$ gives $D_d=d$, as it should at the perpendicular.
Setting $v=0$ gives $\cosh D_d(u,0)=\cosh d\cosh u$:
moving only one endpoint increases the distance. At the
perpendicular, the endpoint Hessian is
\[
 \nabla^2D_d(0,0)=\frac1{\sinh d}
     \begin{pmatrix}\cosh d&-1\\-1&\cosh d\end{pmatrix}.
\]
Its eigenvalues $(\cosh d\pm1)/\sinh d$ are positive. Moving
both endpoints together is cheaper than moving them in opposite
directions, but either motion increases length to second order.

For the ordinary seam between cuffs $a,b>0$ of pants with third cuff
$c\geq0$, the right-angled-hexagon relation gives
\begin{equation}
 \cosh d=\frac{\cosh(c/2)+\cosh(a/2)\cosh(b/2)}
                    {\sinh(a/2)\sinh(b/2)}.
 \label{eq:seam}
\end{equation}
It can also be checked from pants holonomy: translations of lengths
$a,b$ on axes at distance $d$, with the opposed boundary directions,
have product trace
$2\cosh(a/2)\cosh(b/2)-2\cosh d\sinh(a/2)\sinh(b/2)$.
Equating this to $-2\cosh(c/2)$ gives~\eqref{eq:seam}, with the
same lift signs as~\eqref{eq:figuretrace}.
For a returning or winding arc, choose the boundary lifts appropriate to
its relative homotopy class and use their common perpendicular instead.
There is no single seam length that describes all relative arc classes.

For the reduced itinerary of a fixed nonperipheral closed curve, fix the
boundary lifts for each segment and retain its integer winding data.
Cuff gluing then imposes affine relations on the real endpoint positions,
involving the twists $\tau_i$. Minimize the sum of the corresponding
functions~\eqref{eq:localdistance} over the crossing positions. Every
competitor is a piecewise geodesic in the fixed free-homotopy class, so its
length is at least the closed geodesic length. The geodesic representative
realizes equality: geodesics have no bigons with the cutting cuffs and
therefore realize the same reduced itinerary. At the minimum, first
variation imposes smooth tangent matching at every crossing. A curve
contained entirely in one pair of pants is handled directly by its
holonomy; it has no endpoint variables.

The endpoint Hessian will also control variation in a twist. Put
$z=\cosh D_d(u,v)$ and $S=\sqrt{z^2-1}>0$. Direct differentiation gives
\begin{equation}
 \begin{aligned}
 \nabla^2D_d&=\frac{\sinh^2d}{S^3}
 \begin{pmatrix}
 z\cosh^2v&-\cosh u\cosh v\\
 -\cosh u\cosh v&z\cosh^2u
 \end{pmatrix},\\
 \det\nabla^2D_d&=
 \frac{\sinh^4d\cosh^2u\cosh^2v}{S^4}>0.
 \end{aligned}
 \label{eq:localdistancehessian}
\end{equation}
For example, $z_{uu}=z$ and
$z^2-z_u^2-1=\sinh^2d\cosh^2v$, while
$z_{uv}(z^2-1)-zz_uz_v=-\sinh^2d\cosh u\cosh v$.
These identities give the displayed matrix. Its positive diagonal and
determinant prove strict joint convexity in the two endpoint coordinates.
Reversing an endpoint orientation conjugates the Hessian by a sign matrix
and preserves positive definiteness.

\subsection{Holonomy and first variation}

In a basis adapted to a cuff axis, a twist displacement is represented by
\[
 T_i(\tau_i)=\diag(e^{\tau_i/2},e^{-\tau_i/2}).
\]
For a fixed nonperipheral curve word $\gamma$, insert these matrices and
their inverses into the fixed pants
holonomies. The resulting product $W_\gamma$ satisfies
\begin{equation}
 2\cosh(\ell_\gamma/2)
  =\abs{\tr W_\gamma(\mathbf a,T_1(\tau_1),\ldots,T_N(\tau_N))}.
 \label{eq:traceword}
\end{equation}
With the cuff lengths fixed, its trace is a finite Laurent polynomial in
$e^{\tau_i/2}$. Its coefficients need not all be positive in an arbitrary
choice of lifts and coordinates.

The twist derivative has the cosine formula
\begin{equation}
 \frac{\partial\ell_\gamma}{\partial\tau_i}
   =\sum_{p\in\gamma\cap\alpha_i}\cos\vartheta_p,
 \qquad
 \abs{\frac{\partial\ell_\gamma}{\partial\tau_i}}
      \le i(\gamma,\alpha_i).
 \label{eq:cosine}
\end{equation}
Angles and twist orientation are chosen consistently, and intersections
are counted by passages, including repeated passages for an iterate.
Cutting at each crossing proves the cosine formula: endpoint variations
inside a segment cancel by stationarity, while the relative displacement
across the cuff contributes the tangential projection
$\cos\vartheta_p$~\cite{WolpertTwist}. The same argument applies to
immersed geodesics. If $\gamma$ is a power of a cuff, the intersection
sum is empty and the derivative is zero.

Length is strictly convex in one twist $\tau_i$ whenever
$i(\gamma,\alpha_i)>0$. Here is a proof that includes nonsimple curves.
Hold all pants lengths and the other twists fixed, and use the fixed
reduced itinerary above. The segment distances $d$ remain positive and
fixed. For two twist values choose the minimizing endpoint configurations.
Their affine interpolation is admissible for the interpolated twist.
The two full endpoint vectors cannot coincide: at a crossing of
$\alpha_i$, the gluing equation contains $\tau_i$ with coefficient
$1$ or $-1$, with fixed winding data. Thus a change of this twist cannot
be absorbed while leaving all segment endpoints fixed. Strict convexity
of the sum in~\eqref{eq:localdistancehessian}, followed by minimization,
gives the strict convexity inequality for the geodesic length. This also
explains why the conclusion excludes a disjoint cuff, whose twist leaves
the length constant. The simple-curve earthquake theorem provides a
broader context~\cite[Theorem 1]{Kerckhoff}; no extension of that cited
statement is needed for the local argument here.

These ideas overlap with several parts of Mirzakhani's work, but serve
different roles there. Dehn--Thurston coordinates and the definition of
combinatorial length for simple multicurves are given in
\cite[Theorem 3.4 and equation (3.6)]{MirzSimple}. The comparison with
geometric length for an $L$-bounded pants decomposition is
\cite[Proposition 3.5 and equation (3.7), p.~107]{MirzSimple}. The work
on general curve orbits studies asymptotically piecewise-linear length,
including filling curves and length sublevel volumes
\cite[Theorem 4.1 and Section 8]{MirzGeneral}. In the basic volume recursion, the
cutting curves are simple and the kernel depends on their lengths, so the
twist integral itself is just the period $a_i$.

\subsection{Unfolding a twist lattice}

Let $D_i$ be the Dehn twist about the $i$th cuff. With a consistent sign
convention,
\[
 \ell_{D_1^{k_1}\cdots D_N^{k_N}\gamma}(\mathbf a,\boldsymbol\tau)
 =\ell_\gamma(\mathbf a,\boldsymbol\tau+\mathbf k\mathbf a).
\]
For a nonnegative function $F$, Tonelli's theorem gives
\begin{equation}
 \boxed{
 \int_{\prod_i[0,a_i)}\sum_{\mathbf k\in\ZZ^N}
 F\bigl(\ell_{D_1^{k_1}\cdots D_N^{k_N}\gamma}
                  (\mathbf a,\boldsymbol\tau)\bigr)\,\dd\boldsymbol\tau
 =\int_{\RR^N}F\bigl(\ell_\gamma(\mathbf a,\mathbf u)\bigr)\,\dd\mathbf u.}
 \label{eq:twistunfold}
\end{equation}
As written, the sum is over lattice elements, so the nonnegative equality
is valid even if both sides are infinite. For signed functions, impose
absolute integrability on $\RR^N$. To interpret the sum as one term per
curve, first quotient any stabilizer and unfold the corresponding
nonredundant twist directions; twisting about a disjoint cuff must not
produce infinitely many copies of one curve. This quotient may change
the lattice and its fundamental domain, so one cannot simply divide an
infinite multiplicity out of the displayed formula.

A fixed curve's length is not periodic in a twist coordinate. The sum over
its twist orbit is periodic, which is exactly what makes this unfolding
work. The accompanying WP measure is $\prod_i\dd a_i\,\dd\tau_i$.
However, the positive cuff orthant times one twist box is the quotient by
this particular twist subgroup, not a fundamental domain for the full
mapping class group. A further full-moduli integration needs the appropriate
orbit unfolding.

For one cuff, the formula is exactly the circle-to-line toy
unfolding of Section~\ref{sec:unfolding}, with interval length
$a$ instead of $1$ and function $u\mapsto F(\ell_\gamma(u))$.
If the curve is disjoint from that cuff, every twist leaves it
unchanged; this shows concretely why one must remove that
stabilizer direction before counting distinct curves.

For example, with weight $(2\cosh(\ell_\gamma/2))^{-s}$ and
$x_i=e^{u_i/2}$, the right side becomes the Mellin-type integral
\begin{equation}
 2^N\int_{(0,\infty)^N}
 \abs{\tr W_\gamma(\mathbf a,\mathbf x)}^{-s}
             \prod_i\frac{\dd x_i}{x_i},
 \label{eq:mellin}
\end{equation}
whenever it converges. This is an explicit analytic object attached to the
curve itinerary and the fixed pants lengths.

\paragraph{Computational support.}
Groups G4--G5 and G8 of
\path{experiments/quantum_geodesics/README.md} check the local Lorentz
distance, endpoint derivatives and Hessian, seam matrices, and a fixed
stationarity example. They test declared finite inputs, not a general
multicurve minimization or global holonomy-construction algorithm.
The strict-convexity statement rests on the fixed-itinerary argument
above. The result path is \path{artifacts/quantum_geodesics/result.json},
with the formula map in \path{book/claims/sections17-22.json}.
